\clearpage
\begin{explanationbox}
\textbf{一句话直觉}\quad 先把直线限制为经过样本中心，再比较不同倾斜程度的纵向误差；误差平方和最小的那条就是回归线。

% R038：三条候选线采用同一坐标范围、同一单位长度；
% 数据为 (1,2),(2,3),(3,5),(4,6)

\begin{center}
	\begin{tikzpicture}[
			x=1cm,
			y=1cm,
			>=Stealth,
			every node/.style={font=\small},
			axis/.style={->,black!70,thick},
			data/.style={
					circle,
					fill=blue!80!black,
					inner sep=2.3pt
				},
			res/.style={red!78!black,thick}
		]

		\path[use as bounding box] (0,0) rectangle (12.4,7.65);

		\node[
			font=\large\bfseries,
			text=CExplanation
		] at (6.2,7.38)
		{同一坐标尺度下比较：残差平方和最小的候选线最优};

		% 三幅图统一使用相同的横纵坐标缩放
		\def\xstep{0.66}
		\def\ystep{0.50}
		\def\axisY{1.82}

		\foreach \o/\b/\a/\q/\eq/\title in {
		0/1.0/1.5/1.00/{$\hat y=x+1.5$}/{候选},
		4.2/1.4/0.5/0.20/{$\hat y=1.4x+0.5$}/{最优},
		8.4/1.8/-0.5/1.00/{$\hat y=1.8x-0.5$}/{候选}
		}{

		% 卡片背景
		\draw[
			rounded corners=5pt,
			fill=blue!2,
			draw=CExplanation,
			line width=.8pt
		]
		(\o+.08,.48) rectangle (\o+3.92,6.88);

		% 突出中间的最优候选线
		\ifdim\o pt=4.2pt
			\draw[
				orange!85!black,
				line width=1.5pt,
				rounded corners=5pt
			]
			(\o+.08,.48) rectangle (\o+3.92,6.88);
		\fi

		% 卡片标题
		\node[
			font=\bfseries,
			text=CExplanation
		] at (\o+2,6.57)
		{\title：$b=\b$};

		% 回归直线方程
		\node[
			font=\normalsize,
			text=CExplanation
		] at (\o+2,6.08)
		{\eq};

		% x 轴
		\draw[axis]
		(\o+.62,\axisY)
		--
		(\o+3.48,\axisY)
		node[above right=-1pt] {$x$};

		% y 轴
		\draw[axis]
		(\o+.62,\axisY)
		--
		(\o+.62,5.55)
		node[above] {$y$};

		% x 轴刻度
		\foreach \xx in {1,2,3,4}{
				\pgfmathsetmacro{\px}{\o+.62+\xstep*\xx}

				\draw
				(\px,\axisY-.08)
				--
				(\px,\axisY+.08);

				\node[below=3pt] at (\px,\axisY)
				{$\xx$};
			}

		% y 轴刻度
		\foreach \yy in {1,2,3,4,5,6,7}{
				\pgfmathsetmacro{\py}{\axisY+\ystep*\yy}

				\draw
				(\o+.54,\py)
				--
				(\o+.70,\py);

				\node[left=2pt] at (\o+.62,\py)
				{$\yy$};
			}

		% 四个原始数据点
		\foreach \xx/\yy in {1/2,2/3,3/5,4/6}{
				\pgfmathsetmacro{\px}{\o+.62+\xstep*\xx}
				\pgfmathsetmacro{\py}{\axisY+\ystep*\yy}

				\node[data] at (\px,\py) {};
			}

		% 候选回归线的两个端点
		\pgfmathsetmacro{\xone}{\o+.62+\xstep}
		\pgfmathsetmacro{\yone}{\axisY+\ystep*(\a+\b)}

		\pgfmathsetmacro{\xfour}{\o+.62+4*\xstep}
		\pgfmathsetmacro{\yfour}{\axisY+\ystep*(\a+4*\b)}

		% 绘制候选回归线
		\draw[
			CStatement,
			very thick
		]
		(\xone,\yone)
		--
		(\xfour,\yfour);

		% 绘制纵向残差
		\foreach \xx/\yy in {1/2,2/3,3/5,4/6}{
				\pgfmathsetmacro{\px}{\o+.62+\xstep*\xx}
				\pgfmathsetmacro{\py}{\axisY+\ystep*\yy}
				\pgfmathsetmacro{\ly}{\axisY+\ystep*(\a+\b*\xx)}

				\draw[res]
				(\px,\py)
				--
				(\px,\ly);
			}

		% 样本中心点
		\pgfmathsetmacro{\cx}{\o+.62+2.5*\xstep}
		\pgfmathsetmacro{\cy}{\axisY+4*\ystep}

		\fill[orange!90!black]
		(\cx,\cy) circle (3.1pt);

		\node[
			anchor=south,
			xshift=-10pt,
			text=orange!85!black
		] at (\cx,\cy)
		{$(2.5,4)$};

		% 残差平方和：单独留出底部空间
		\node[
			font=\bfseries\small,
			text=red!78!black
		] at (\o+2,.91)
		{$Q=\sum e_i^2=\q$};
		}

		% 图下注释
		\node[
			text=red!78!black,
			anchor=west,
			font=\small
		] at (.22,.14)
		{每条红线段均为纵向残差；只画三条候选线，严格全局最优见后文。};

	\end{tikzpicture}
\end{center}


上图替代“反复试斜率、逐点计算”的长叙述：三个候选线都过 $(2.5,4)$，但只有 $\hat y=1.4x+0.5$ 使 $Q=0.20$ 最小。三个小图使用相同的横、纵坐标范围和单位长度。

\clearpage
% R038：纵向残差与最短垂距严格区分；
% 在 x=1cm、y=1.18cm 的画布尺度下，灰线与回归线严格垂直。

\begin{center}
	\begin{tikzpicture}[
			x=1cm,
			y=1.18cm,
			>=Stealth,
			every node/.style={font=\small}
		]

		\path[use as bounding box] (0,0) rectangle (12.4,5.8);

		\node[font=\large\bfseries,text=CExplanation] at (6.2,5.46)
		{最小二乘最小化纵向残差平方，不是点到直线的最短距离};

		% -------------------- 左侧：两种距离 --------------------
		\draw[rounded corners=5pt,fill=blue!2,draw=CExplanation]
		(.15,.28) rectangle (6.02,4.98);

		\draw[->,black!70] (.72,.82)--(5.56,.82) node[right] {$x$};
		\draw[->,black!70] (.72,.82)--(.72,4.35) node[above] {$y$};

		% 回归线
		\draw[CStatement,very thick] (.92,1.25)--(5.28,3.90);

		% 数据点
		\fill[blue!80!black] (3.65,3.88) circle (3.1pt)
		node[above=4pt] {数据点 $(x_i,y_i)$};

		% 红线：固定 x_i 后的纵向残差
		\draw[red!80!black,very thick]
		(3.65,3.88)--(3.65,2.91);

		\node[
			anchor=east,
			text=red!80!black,
			font=\scriptsize
		] at (3.50,3.36)
		{$e_i=y_i-\hat y_i$};

		% 灰线：点到回归线的最短垂距
		\coordinate (foot) at (4.1925,3.2390);

		\draw[gray!80,very thick]
		(3.65,3.88)--(foot);

		\node[
			anchor=west,
			text=gray!80,
			font=\scriptsize
		] at (4.08,2.90)
		{灰：最短垂距};

		% 直角标记
		\draw[gray!80,thick]
		(4.10,3.35)--(3.97,3.27)--(4.06,3.16);

		\node[text=CExplanation,font=\bfseries] at (3.10,.54)
		{红：固定 $x_i$ 的纵向差};

		% -------------------- 右侧：为什么平方 --------------------
		\draw[rounded corners=5pt,fill=orange!3,draw=orange!80!black]
		(6.40,.28) rectangle (12.25,4.98);

		\node[font=\bfseries,text=orange!85!black] at (9.32,4.48)
		{为什么要平方？};

		% 正负误差分开显示
		\draw[red!75!black,->,thick]
		(7.90,3.64)--(9.05,3.64)
		node[midway,above=2pt] {$+2$};

		\draw[red!75!black,->,thick]
		(10.74,3.64)--(9.59,3.64)
		node[midway,above=2pt] {$-2$};

		\node at (9.32,3.08)
		{$+2+(-2)=0$：会抵消};

		% 三个正方形：边长按 1:2:3，面积按 1:4:9
		% y 单位为 1.18cm，因此纵坐标高度除以 1.18

		\draw[fill=orange!25,draw=orange!85!black]
		(7.15,1.42) rectangle (7.65,1.844);
		\node at (7.40,1.632) {$1$};

		\draw[fill=orange!25,draw=orange!85!black]
		(8.82,1.42) rectangle (9.82,2.267);
		\node at (9.32,1.844) {$4$};

		\draw[fill=orange!25,draw=orange!85!black]
		(10.40,1.42) rectangle (11.90,2.691);
		\node at (11.15,2.056) {$9$};

		\node at (7.40,1.07) {$1^2=1$};
		\node at (9.32,1.07) {$2^2=4$};
		\node at (11.15,1.07) {$3^2=9$};

		\node[font=\small\bfseries] at (9.32,.58)
		{不抵消，且大误差惩罚更重};

	\end{tikzpicture}
\end{center}

\textbf{读图关键}\quad 普通最小二乘固定 $x_i$，因此衡量的是 $y_i-\hat y_i$ 的平方；它不是点到直线的最短垂距平方。平方既避免正负误差抵消，也更严厉地惩罚大误差。

\clearpage
% R038：样本中心与完整四象限的离差乘积符号图。
\begin{center}
	\begin{tikzpicture}[
			x=1cm,
			y=1.15cm,
			>=Stealth,
			every node/.style={font=\small}
		]

		\path[use as bounding box] (0,0) rectangle (12.4,7.0);

		\node[font=\large\bfseries,text=CExplanation]
		at (6.2,6.66)
		{样本中心固定位置；离差乘积和决定斜率方向};

		%==================== 左侧：候选直线 ====================

		\draw[rounded corners=5pt,fill=blue!2,draw=CExplanation]
		(.15,.28) rectangle (6.02,6.12);

		\node[
			font=\bfseries\scriptsize,
			text=CExplanation,
			text width=4.8cm,
			align=center,
			inner sep=0pt
		] at (3.08,5.67)
		{确定斜率后，截距由样本中心确定};

		% 坐标图整体上移，给底部说明留出独立空间
		\draw[->,black!65]
		(.70,1.35)--(5.45,1.35) node[right] {$x$};

		\draw[->,black!65]
		(.70,1.35)--(.70,5.25) node[above] {$y$};

		% 三条直线都经过样本中心
		\draw[gray!70,thick]
		(.95,3.776)--(4.75,2.216);

		\draw[gray!70,thick]
		(.87,2.488)--(4.83,3.568);

		\draw[CStatement,very thick]
		(1.10,1.780)--(4.70,4.390);

		% 样本点
		\foreach \x/\y in {
				1.12/1.82,
				1.70/2.22,
				2.48/2.96,
				3.12/2.82,
				3.80/3.92,
				4.60/4.35
			}{
				\fill[blue!80!black] (\x,\y) circle (2.5pt);
			}

		\coordinate (c) at (2.8033,3.0150);

		\fill[orange!90!black] (c) circle (4pt);

		\node[
			font=\scriptsize,
			text=orange!85!black,
			anchor=south,
			fill=blue!2,
			inner sep=1pt
		] at (2.8033,3.19)
		{$(\bar x,\bar y)$};

		\draw[->,orange!85!black,thick]
		(4.32,1.88) arc (-42:35:1.45);

		\node[
			font=\scriptsize,
			text=CExplanation,
			text width=4.9cm,
			align=center,
			inner sep=0pt
		] at (3.08,.84)
		{$\hat a=\bar y-\hat b\bar x$\\[-1pt]
			候选线均绕样本中心转动};

		%==================== 右侧：四象限符号 ====================

		\draw[rounded corners=5pt,fill=green!2,draw=green!55!black]
		(6.40,.28) rectangle (12.25,6.12);

		\node[
			font=\bfseries\scriptsize,
			text=green!45!black,
			text width=5.2cm,
			align=center,
			inner sep=0pt
		] at (9.32,5.73)
		{四象限：每个矩形对应一个离差点};

		% 四象限整体上移
		\coordinate (O) at (9.35,3.68);

		\draw[densely dashed,gray]
		(6.82,3.68)--(11.82,3.68);

		\draw[densely dashed,gray]
		(9.35,1.78)--(9.35,5.35);

		\fill[orange!90!black] (O) circle (3.5pt);

		\node[font=\scriptsize,above right,inner sep=1.5pt]
		at (O) {中心};

		% 四个带符号矩形
		\fill[green!35] (O) rectangle (10.75,4.95);
		\draw[black!60] (O) rectangle (10.75,4.95);

		\fill[red!25] (O) rectangle (7.92,4.68);
		\draw[black!60] (O) rectangle (7.92,4.68);

		\fill[green!35] (O) rectangle (8.08,2.22);
		\draw[black!60] (O) rectangle (8.08,2.22);

		\fill[red!25] (O) rectangle (10.82,2.15);
		\draw[black!60] (O) rectangle (10.82,2.15);

		% 四个离差点
		\fill[blue!80!black] (10.75,4.95) circle (2.7pt);
		\fill[blue!80!black] (7.92,4.68) circle (2.7pt);
		\fill[blue!80!black] (8.08,2.22) circle (2.7pt);
		\fill[blue!80!black] (10.82,2.15) circle (2.7pt);

		% 点名全部放进矩形内部
		\node[font=\scriptsize,below left=3pt]
		at (10.75,4.95) {$P_1$};

		\node[font=\scriptsize,below right=3pt]
		at (7.92,4.68) {$P_2$};

		\node[font=\scriptsize,above right=3pt]
		at (8.08,2.22) {$P_3$};

		\node[font=\scriptsize,above left=3pt]
		at (10.82,2.15) {$P_4$};

		% 象限符号
		\node[font=\bfseries] at (10.05,4.315) {$+$};
		\node[font=\bfseries] at (8.635,4.18) {$-$};
		\node[font=\bfseries] at (8.715,2.95) {$+$};
		\node[font=\bfseries] at (10.085,2.915) {$-$};

		% 底部说明与四象限完全分离
		\node[
			font=\scriptsize,
			align=center,
			text=green!45!black,
			text width=5.3cm,
			inner sep=0pt
		] at (9.32,.99)
		{带符号面积：$S_i=(x_i-\bar x)(y_i-\bar y)$\\[-1pt]
			几何面积：$|S_i|=|x_i-\bar x|\,|y_i-\bar y|$\\[-1pt]
			正负由象限决定\\[-1pt]
			“正面积占优”：$\sum_i S_i>0$（代数和为正）};

	\end{tikzpicture}
\end{center}

分子 $S_{xy}=\sum(x_i-\bar x)(y_i-\bar y)$ 可看作带符号面积之和；分母 $S_{xx}$ 衡量 $x$ 的横向展开。故 $S_{xy}$ 的正负决定斜率方向，且 $S_{xx}>0$ 才能唯一估计斜率。
\end{explanationbox}
